\scanpage{112}
\margin{\diagram{\begin{tikzcd}[row sep=small,column sep=small]
1\arrow[r,dash] & \unsure{\mathbb Z/2\mathbb Z}\arrow[r,dash] & \Gamma\arrow[r,dash] & \Gamma'\arrow[r,dash] & 0 \\
& 1\arrow[u,dash]\arrow[r,dash] & \Gamma_{\mathfrak p}\arrow[u,hook]
\end{tikzcd}}}
D'autre part, $(\mathbb Z^P)'$ correspond à un tore sur $Q$, qui peut se décrire ainsi, comme on voit par th. de Galois: l'ext. quadratique $L_0$ de $K_0$ définit une forme tordue $U$ de $G_m$ (savoir $\operatorname{Ker}(\prod_{L_0/K_0}G_{mL_0}\overset{\text{Norm}}{\unsure{\text{—}|\longrightarrow}}G_{mK_0})$) et on considère le
\[
G'=\prod_{K_0/Q}U=\operatorname{Ker}\bigl(\prod_{L_0/Q}G_{mL_0}\mathbin{-}\prod_{K_0/Q}G_{mK_0}\bigr)
\]
(donc $G'(Q)=\operatorname{Ker}(L_0^*\overset{N_{L_0/K_0}}{\unsure{\text{—}|\longrightarrow}}K_0^*)$). On a donc un hom
\diagram{\begin{tikzcd}
(\unsure{**})\quad\varphi:G'\arrow[r] & G
\end{tikzcd}}
qui est une isogénie \uline{de $G'$ sur $G^{0u}$}.

D'ailleurs $\nu_{\unsure{L_0}}:G_{mL_0}\mathbin{-}G_{L_0}$ se factorise par $G'_{L_0}$
\diagram{\begin{tikzcd}
\unsure{\text{via}}\quad G_{mL_0}\arrow[r,dash,"\nu'"]\arrow[d,equal] & G'_{L_0}\arrow[d] \\
G_{mL_0}\arrow[r,dash,"\nu"] & G_{L_0}
\end{tikzcd}}

On fera attention que l'élém $f_0$ \inserted{$=f\,j(p^{1/2})$} (de $G^{0u}(Q)$) ne se relève en un élém \deleted{\unsure{de}} $f'$ de $G'(Q)$ que si \deleted{\unsure{q}}\inserted{$\varphi$} est un isom i.e. $M(q)/\text{torsion}\xrightarrow{\sim}(\mathbb Z^P)'$.
\mspage{111}
